\chapter{Property Cursor Class Hierarchy}
\label{chap:appendix:prop-cursors}

This appendix provides the class hierarchy and annotated code listings for the property cursor abstraction introduced in \autoref{arch:prop-cursors}.
\autoref{fig:prop-cursor-hierarchy} shows the inheritance diagram, and the listings that follow illustrate how the same \code{next()} call leads to different \wt{} operations depending on the storage mode.

%%%%%%%%%%%%%%%%%%%%%%%%%%%%%%%%%%%%%%%%%%%%%%%%%%%%%%%%%%%%%%%%%%%%%%%%%%%%%%%%
\section{Class Hierarchy}
\label{appendix:prop-cursor-hierarchy}

\autoref{fig:prop-cursor-hierarchy} shows the property cursor inheritance hierarchy.
The two abstract base classes, \code{NodePropCursor} and \code{EdgePropCursor}, define the mode-agnostic interface.
Each storage mode provides a pair of concrete implementations.
Columnar mode uses \code{WTNodePropCursor} and \code{WTEdgePropCursor}.
Embedded mode uses \code{EmbNodePropCursor} and \code{EmbEdgePropCursor}.

\begin{figure}[htbp]
  \centering
  \begin{tikzpicture}[
    class/.style={rectangle, draw, minimum width=3.0cm, minimum height=0.8cm, font=\small\ttfamily},
    abstract/.style={class, dashed},
    arr/.style={-{Triangle[length=8pt, width=6pt]}, thick}
  ]
    % NodePropCursor hierarchy
    \node[abstract] (npc) at (0, 0) {NodePropCursor};
    \node[class] (wtnpc) at (-3.5, -1.8) {WTNodePropCursor};
    \node[class] (embnpc) at (3.5, -1.8) {EmbNodePropCursor};

    % EdgePropCursor hierarchy (stacked below to avoid overlap)
    \node[abstract] (epc) at (0, -4.5) {EdgePropCursor};
    \node[class] (wtepc) at (-3.5, -6.3) {WTEdgePropCursor};
    \node[class] (embepc) at (3.5, -6.3) {EmbEdgePropCursor};

    % Inheritance arrows
    \draw[arr] (wtnpc) -- (npc);
    \draw[arr] (embnpc) -- (npc);
    \draw[arr] (wtepc) -- (epc);
    \draw[arr] (embepc) -- (epc);

    % Labels
    \node[font=\scriptsize, below=0.1cm] at (wtnpc.south) {(columnar)};
    \node[font=\scriptsize, below=0.1cm] at (embnpc.south) {(embedded)};
    \node[font=\scriptsize, below=0.1cm] at (wtepc.south) {(columnar)};
    \node[font=\scriptsize, below=0.1cm] at (embepc.south) {(embedded)};
  \end{tikzpicture}
  \caption{Property cursor class hierarchy.  Dashed boxes are abstract base classes; solid boxes are concrete implementations selected by the \code{GraphBase} factory methods based on the storage mode.}
  \label{fig:prop-cursor-hierarchy}
\end{figure}

%%%%%%%%%%%%%%%%%%%%%%%%%%%%%%%%%%%%%%%%%%%%%%%%%%%%%%%%%%%%%%%%%%%%%%%%%%%%%%%%
\section{Abstract Interfaces}
\label{appendix:prop-cursor-interfaces}

\autoref{lst:node-prop-cursor} shows the \code{NodePropCursor} abstract interface.
\code{EdgePropCursor} follows the same pattern, replacing vertex-range positioning with source-vertex pinning.

\begin{lstlisting}[style=cpp, label={lst:node-prop-cursor}, caption={\code{NodePropCursor} abstract interface (from \texttt{iterator.h}).}, captionpos=b]
class NodePropCursor {
public:
    virtual ~NodePropCursor() = default;

    // Positioning
    virtual void set_range(node_id_t start,
                           node_id_t end) = 0;
    virtual void seek(node_id_t id) = 0;

    // Iteration
    virtual bool next() = 0;
    virtual node_id_t current_id() const = 0;

    // Typed accessors (column index n)
    virtual uint64_t get_uint64(int n) const = 0;
    virtual int64_t  get_int64(int n) const = 0;
    virtual int32_t  get_int32(int n) const = 0;
    virtual std::string_view get_string(int n)
                                       const = 0;

    virtual void close() = 0;
};
\end{lstlisting}

\begin{lstlisting}[style=cpp, label={lst:edge-prop-cursor}, caption={\code{EdgePropCursor} abstract interface (from \texttt{iterator.h}).}, captionpos=b]
class EdgePropCursor {
public:
    virtual ~EdgePropCursor() = default;

    // Pin a single source vertex
    virtual void set_src(node_id_t src) = 0;
    // Bulk scan over a source range
    virtual void set_src_range(node_id_t start,
                               node_id_t end) = 0;

    virtual bool next() = 0;
    virtual node_id_t current_src() const = 0;
    virtual node_id_t current_dst() const = 0;

    virtual uint64_t get_uint64(int n) const = 0;
    virtual int64_t  get_int64(int n) const = 0;
    virtual int32_t  get_int32(int n) const = 0;
    virtual std::string_view get_string(int n)
                                       const = 0;

    virtual void close() = 0;
};
\end{lstlisting}

%%%%%%%%%%%%%%%%%%%%%%%%%%%%%%%%%%%%%%%%%%%%%%%%%%%%%%%%%%%%%%%%%%%%%%%%%%%%%%%%
\section{Column-Group Layout Dispatch}
\label{appendix:colgroup-layout}

Each column group has a value format that determines the arity and types of the \code{get\_value()} call to \wt{}.
\autoref{lst:colgroup-layout} shows the \code{ColgroupValueLayout} enum and the dispatch logic.
The layout avoids unpacking columns that belong to a different group: if a column group stores two \code{uint64} values, the cursor calls \code{get\_value(cursor, \&v1, \&v2)} rather than unpacking the full record.

\begin{lstlisting}[style=cpp, label={lst:colgroup-layout}, caption={\code{ColgroupValueLayout} enum and accessor dispatch.}, captionpos=b]
enum ColgroupValueLayout {
    CVL_Q,     // single uint64
    CVL_QQ,    // two uint64s
    CVL_Qi,    // uint64 + int32
    CVL_SS,    // two strings
    CVL_SSS,   // three strings
    // ... additional layouts as needed
};

ColgroupValueLayout layout_for(
    const std::string& colgroup_name,
    int label_id);
\end{lstlisting}

%%%%%%%%%%%%%%%%%%%%%%%%%%%%%%%%%%%%%%%%%%%%%%%%%%%%%%%%%%%%%%%%%%%%%%%%%%%%%%%%
\section{Columnar vs.\ Embedded \texttt{next()} Paths}
\label{appendix:next-paths}

The two storage modes differ in how \code{next()} retrieves the property record.
\autoref{lst:wt-next} shows the columnar path: the cursor advances the \wt{} column-group cursor and calls \code{get\_value()} with an arity matching the column group's layout.
\autoref{lst:emb-next} shows the embedded path: the cursor reads the raw property blob from the structural table (or the separate \code{node\_props} table in AdjList mode) and applies extractor function pointers to pull typed fields from fixed byte offsets.

\begin{lstlisting}[style=cpp, label={lst:wt-next}, caption={Columnar path: \code{WTNodePropCursor::next()}.}, captionpos=b]
bool WTNodePropCursor::next() {
    int ret = wt_cursor->next(wt_cursor);
    if (ret != 0) return false;

    // Read key (big-endian vertex ID)
    WT_ITEM key_item;
    wt_cursor->get_key(wt_cursor, &key_item);
    cur_id = decode_be64(key_item.data);

    if (cur_id >= range_end) return false;

    // Read value: arity depends on layout
    switch (layout) {
    case CVL_QQ:
        wt_cursor->get_value(wt_cursor,
                             &vals_u64[0],
                             &vals_u64[1]);
        break;
    case CVL_Q:
        wt_cursor->get_value(wt_cursor,
                             &vals_u64[0]);
        break;
    // ... other layouts
    }
    return true;
}
\end{lstlisting}

\begin{lstlisting}[style=cpp, label={lst:emb-next}, caption={Embedded path: \code{EmbNodePropCursor::next()}.}, captionpos=b]
bool EmbNodePropCursor::next() {
    int ret = wt_cursor->next(wt_cursor);
    if (ret != 0) return false;

    WT_ITEM key_item, val_item;
    wt_cursor->get_key(wt_cursor, &key_item);
    cur_id = decode_be64(key_item.data);

    if (cur_id >= range_end) return false;

    // Read raw property blob
    wt_cursor->get_value(wt_cursor, &val_item);
    prop_blob = static_cast<const uint8_t*>(
                    val_item.data);

    // Typed access via extractor pointers:
    //   get_uint64(n) returns
    //     extractors_u64[n](prop_blob)
    // Each extractor reads from a fixed
    // byte offset in the blob.
    return true;
}
\end{lstlisting}

The key difference is in how data reaches the application.
In the columnar path, \wt{} itself unpacks the column group's typed format, and only pages belonging to that column group are loaded into the buffer pool.
In the embedded path, \wt{} delivers an opaque byte array, and the cursor applies schema-specific extractors to interpret it.
The columnar path trades one layer of indirection for the I/O savings of reading fewer pages.
